\documentclass[11pt]{amsart}
\usepackage{amsmath,amssymb,amsthm}
\usepackage[margin=1.2in]{geometry}
\usepackage{booktabs}
\usepackage[colorlinks=true,linkcolor=blue,citecolor=blue]{hyperref}

\newtheorem{theorem}{Theorem}[section]
\newtheorem{lemma}[theorem]{Lemma}
\newtheorem{corollary}[theorem]{Corollary}
\newtheorem{proposition}[theorem]{Proposition}
\theoremstyle{definition}
\newtheorem{definition}[theorem]{Definition}
\newtheorem{example}[theorem]{Example}
\theoremstyle{remark}
\newtheorem{remark}[theorem]{Remark}

\newcommand{\Q}{\mathbb{Q}}
\newcommand{\Z}{\mathbb{Z}}
\newcommand{\R}{\mathbb{R}}
\newcommand{\C}{\mathbb{C}}
\newcommand{\D}{\mathbb{D}}
\newcommand{\lcm}[1]{[1,\dots,#1]}

\title{A rank-sensitive strengthening of the determinantal\\
irrationality criterion for Ap\'ery limits}
\date{June 2026 --- draft for expert scrutiny}

\begin{document}

\begin{abstract}
Let $A(x)=\sum u_nx^n$ and $V(x)=\sum v_nx^n$ be annihilated by a common
differential operator $L$ of order $k\ge 2$, irreducible over $\Q(x)$, with
$u_n\in\Z$ and $c\,\lcm{n}^{\sigma}v_n\in\Z$. Suppose that for some real
$\eta$ the linear-form generating function $\eta A-V\not\equiv 0$ extends
holomorphically to the slit plane $\C\setminus[x_2,\infty)$. We observe
that the arithmetic holonomy bound of Calegari, Dimitrov and Tang, applied
to the tuple $\{1,P,\theta P,\dots,\theta^{k-1}P\}$ with
$\theta=x\,\frac{d}{dx}$, yields: if
\[
\log(4x_2) \;>\; \sigma\,\frac{k+2}{k+1},
\]
then $\eta\notin\Q$. The case $k=1$-analogue
$\log(4x_2)>\tfrac32\sigma$ is Zudilin's determinantal criterion, and the
classical Ap\'ery-type criterion requires $\log x_2>\sigma$; the present
inequality is strictly weaker than both for every $k\ge2$, and opens the
window $e^{\sigma(k+2)/(k+1)}/4 < x_2 \le e^{\sigma}$ in which no previous
criterion applies. We tabulate the resulting specification curve for
linear forms in $1$ and $\zeta(5)$, and we report numerical experiments
indicating that the denominators realized by the underlying auxiliary
constructions are substantially smaller than the theoretical bound, with
the savings concentrated at small primes.
\end{abstract}

\maketitle

\section{Introduction}

Let $(u_n)_{n\ge0}\subset\Z$ and $(v_n)_{n\ge0}\subset\Q$, and let
$\eta\in\R$ be a target number approximated by the linear forms
$\ell_n:=u_n\eta-v_n$. Write $\lcm{n}$ for $\mathrm{lcm}(1,\dots,n)$ and
suppose
\begin{equation}\label{eq:type}
c\,\lcm{n}^{\sigma}\,v_n\in\Z \qquad (n\ge 0,\ \text{with the convention }\lcm{0}:=1)
\end{equation}
for fixed positive integers $c$ and $\sigma$; the condition at $n=0$,
namely $cv_0\in\Z$, may always be arranged by enlarging $c$, since
$v_0$ is a single rational number. Ap\'ery's method
\cite{Apery,vdP} proves $\eta\notin\Q$ as soon as
$\limsup|\ell_n|^{1/n}=:1/x_2$ satisfies $x_2>e^{\sigma}$: indeed, if
$\eta=p/q$ then $cq\lcm{n}^{\sigma}\ell_n$ are integers of absolute value
$\to 0$, hence eventually zero.

Zudilin \cite{Zudilin} extended the reach of the method to situations
where the \emph{decay} of $\ell_n$ is insufficient but the generating
function
\[
L_\eta(x):=\sum_{n\ge0}\ell_n x^n=\eta A(x)-V(x),\qquad
A=\sum u_nx^n,\; V=\sum v_nx^n,
\]
extends holomorphically to the slit plane
$\Omega:=\C\setminus[x_2,\infty)$: under a positivity hypothesis his
determinantal criterion proves $\eta\notin\Q$ whenever
$\log(4x_2)>\tfrac32\sigma$. The constant $4x_2$ is the conformal radius
of $\Omega$ at the origin, and the exponent $\tfrac32$ reflects the
arithmetic cost of Hankel determinants. The same inequality reappears,
with the positivity hypothesis removed, as the $m=2$ case of the
\emph{univalent arithmetic holonomy bound} of
Calegari--Dimitrov--Tang \cite[Theorem~2.7.1 and \S2.7.7]{CDT}.

The purpose of this note is to record the observation that the
\emph{order of the differential equation} satisfied by the linear forms
is an arithmetic resource: feeding the derivative tower of the
(conditional) rational function $P=pA-qV$ into the holonomy bound of
\cite{CDT} lowers the exponent $\tfrac32$ to $\tfrac{k+2}{k+1}$, where
$k$ is the order of the underlying irreducible operator.

\begin{theorem}\label{thm:main}
Let $k\ge2$ and let $L$ be a linear differential operator of order $k$
with coefficients in $\Q(x)$, irreducible over $\Q(x)$. Let
$A=\sum u_nx^n$ and $V=\sum v_nx^n$ in $\Q[[x]]$ be annihilated by $L$,
with $u_n\in\Z$ and with \eqref{eq:type} for fixed positive integers
$c,\sigma$. Let $\eta\in\R$ be such that $L_\eta=\eta A-V$ is not
identically zero and extends holomorphically to
$\Omega=\C\setminus[x_2,\infty)$ for some $x_2>0$. If
\begin{equation}\label{eq:criterion}
\log(4x_2)\;>\;\sigma\,\frac{k+2}{k+1},
\end{equation}
then $\eta\notin\Q$.
\end{theorem}

For $\sigma=3$ (the type of Ap\'ery's $\zeta(3)$ approximations) the
thresholds on $x_2$ read:
\begin{center}
\begin{tabular}{lcc}
\toprule
criterion & condition & threshold on $x_2$\\
\midrule
classical (radius) & $x_2>e^{\sigma}$ & $20.09$\\
determinantal \cite{Zudilin} & $4x_2>e^{3\sigma/2}$ & $22.51$\\
Theorem~\ref{thm:main}, $k=2$ & $4x_2>e^{4\sigma/3}$ & $13.65$\\
Theorem~\ref{thm:main}, $k=3$ & $4x_2>e^{5\sigma/4}$ & $10.63$\\
Theorem~\ref{thm:main}, $k\to\infty$ & $4x_2>e^{\sigma}$ & $5.02$\\
\bottomrule
\end{tabular}
\end{center}
Every $k\ge2$ row opens territory where neither the classical nor the
determinantal criterion concludes. We stress at once that for Ap\'ery's
actual $\zeta(3)$ data ($x_2=(1+\sqrt2)^4\approx 33.97$, $k=3$) the
classical criterion already applies; the content of
Theorem~\ref{thm:main} lies entirely in the window
$x_2\in\bigl(e^{\sigma(k+2)/(k+1)}/4,\;e^{\sigma}\bigr]$, i.e.\ in its
potential applications to constructions of sub-classical quality, for
instance towards $\zeta(5)$ (\S\ref{sec:window}).

Theorem~\ref{thm:main} is a direct application of the published
holonomy bound of \cite{CDT}; the proof below consists of the
verification, which we have endeavoured to make complete, that the
derivative tower satisfies every hypothesis of that bound, together
with the elementary computation of the resulting numerology. Given that
the consequence strengthens two classical criteria on a well-trodden
battlefield, we submit this note explicitly for expert scrutiny.

\section{The input from \cite{CDT}}

We quote the univalent arithmetic holonomy bound in the form we use.
For an $m\times r$ array $\mathbf b=(b_{i,j})$ of nonnegative reals
whose columns consist of $u_j$ zeros followed by a constant value
$b_j$, set $\sigma_i:=\sum_j b_{i,j}$ (row sums, nondecreasing) and
\[
\tau(\mathbf b):=\frac1{m^2}\sum_{i=1}^m(2i-1)\,\sigma_i.
\]

\begin{theorem}[{\cite[Theorem~2.7.1; see also Theorem~7.0.1 and the
holonomic relaxation stated there]{CDT}}]\label{thm:cdt}
Let $\Omega\subset\C$ be a contractible domain containing $0$, of
conformal radius $\rho=\rho(\Omega,0)$. Let $f_1,\dots,f_m\in\Q[[x]]$ be
$\Q(x)$-linearly independent, holonomic, meromorphic on $\Omega$, with
\[
f_i(x)=a_{i,0}+\sum_{n\ge1}a_{i,n}\,
\frac{x^n}{\lcm{b_{i,1}n}\cdots\lcm{b_{i,r}n}},\qquad a_{i,n}\in\Z,
\]
for the array $\mathbf b$ as above, and suppose
$\rho>e^{\tau(\mathbf b)}$. Then
\[
m\;\le\;\frac{\log\rho}{\log\rho-\tau(\mathbf b)}.
\]
\end{theorem}

We will also use the classical fact that the Koebe-type map
$\varphi(z)=4x_2\,z/(1+z)^2$ is a conformal equivalence
$\D\to\C\setminus[x_2,\infty)$ with $\varphi(0)=0$ and
$\varphi'(0)=4x_2$; thus $\rho(\Omega,0)=4x_2$ for the slit plane.

\begin{remark}\label{rem:shidlovsky}
The vanishing-filtration input behind Theorem~\ref{thm:cdt} is, for the
tuple used below, available in its oldest form: the $\Q(x)$-linear span
of $\{1,P,\theta P,\dots,\theta^{k-1}P\}$ is closed under $d/dx$
(see Lemma~\ref{lem:types}), so Shidlovsky's lemma
\cite[Theorem~3.2.8]{CDT} applies with no $\varepsilon$-loss, and the
system is Fuchsian, so Chudnovsky's effective form
\cite[Theorem~3.2.10, Remark~3.2.12]{CDT} applies as well.
\end{remark}

\section{Proof of Theorem~\ref{thm:main}}

Assume $\eta=p/q$ with $p,q$ coprime integers, $q\ge1$. Set
\[
P:=q\,L_\eta=pA-qV\in\Q[[x]],\qquad \theta:=x\frac{d}{dx}.
\]
Then $P\not\equiv0$, $L(P)=0$, and $P$ is holomorphic on
$\Omega=\C\setminus[x_2,\infty)$.

\begin{lemma}[Independence]\label{lem:indep}
Let $L$ be irreducible over $\Q(x)$ of order $k\ge2$ and let
$0\ne P\in\Q[[x]]$ satisfy $L(P)=0$. Then
$1,P,\theta P,\dots,\theta^{k-1}P$ are $\Q(x)$-linearly independent.
\end{lemma}

\begin{proof}
Since $x$ is invertible in $\Q(x)$, the span of
$\{\theta^jP\}_{j<k}$ over $\Q(x)$ equals that of
$\{P,P',\dots,P^{(k-1)}\}$; suppose
$r_0+M(P)=0$ with $r_0\in\Q(x)$ and $M$ an operator of order $\le k-1$
over $\Q(x)$, not both zero. In the left-Euclidean ring
$\Q(x)\langle\partial\rangle$ the annihilator of $P$ is the left ideal
generated by the minimal operator $L_P$ of $P$, and every annihilator
of $P$ is a left multiple of $L_P$. As $L$ annihilates $P$ and is
irreducible, $L_P$ has order $0$ or $k$; order $0$ forces $P=0$, so
$L_P$ has order $k$.

If $r_0=0$ then $M\ne0$ annihilates $P$ with order $\le k-1<k$, a
contradiction. If $r_0\ne0$ then $M\ne 0$ (else $r_0=0$), and from
$M(P)=-r_0$ we get
$\bigl(\partial-\tfrac{r_0'}{r_0}\bigr)M\,(P)=-r_0'+\tfrac{r_0'}{r_0}r_0=0$.
Thus $N:=(\partial-\tfrac{r_0'}{r_0})M$ is a nonzero annihilator of $P$
of order $\le k$, hence a left multiple of $L_P$:
$N=c\,L_P$ with $c\in\Q(x)^{\times}$ by comparing orders, and therefore
$L_P=c^{-1}(\partial-\tfrac{r_0'}{r_0})\,M$ admits the right-hand
factor $M$ of order $k-1$ with $0<k-1<k$, contradicting irreducibility
(note $\mathrm{ord}\,M=k-1$ is forced by $\mathrm{ord}\,N=k$; and
$\mathrm{ord}\,M=0$ would make $P\in\Q(x)$, i.e.\ $L_P$ of order
$1<k$).
\end{proof}

\begin{lemma}[Types and closure]\label{lem:types}
For $0\le j\le k-1$, the function $cq\,\theta^jP$ has coefficients
$cq\,n^j(pu_n-qv_n)$, and
$\lcm{n}^{\sigma}\cdot cq\,n^j(pu_n-qv_n)\in\Z$; thus each $\theta^jP$
is, after multiplication by the fixed integer $cq$, of the arithmetic
type of Theorem~\ref{thm:cdt} with row $(1,\dots,1)$ ($\sigma$ columns,
$b_j=1$). Moreover the $\Q(x)$-span of
$\{1,P,\theta P,\dots,\theta^{k-1}P\}$ is closed under $d/dx$.
\end{lemma}

\begin{proof}
The coefficient statement is immediate from \eqref{eq:type} and
$u_n\in\Z$; multiplication of a coefficient by the integer $n^j$
preserves integrality against the same denominators. For closure:
multiplying $L$ by $x^k$ and rewriting in $\theta$ (using
$x^i\partial^i\in\Z[\theta]$) expresses $\theta^kP$ as a
$\Q(x)$-combination of $P,\dots,\theta^{k-1}P$; and
$\frac{d}{dx}\theta^jP=x^{-1}\theta^{j+1}P$ stays in the span, while
$\frac{d}{dx}1=0$.
\end{proof}

\begin{lemma}[Local normalization at no arithmetic cost]\label{lem:val}
Write $P=\sum_n c_nx^n$, set $h_0:=P-P(0)\cdot 1$ and, for
$1\le j\le k-1$,
\[
h_j\;:=\;\prod_{i=1}^{j}(\theta-i)\,P\;-\;(-1)^j j!\,P(0)\cdot 1 .
\]
Then $\{1,h_0,h_1,\dots,h_{k-1}\}$ is obtained from
$\{1,P,\theta P,\dots,\theta^{k-1}P\}$ by an integral, triangular,
unit-diagonal (hence unimodular) change of basis; each $cq\,h_j$
carries the same arithmetic type (row $(1,\dots,1)$ of length
$\sigma$) as $cq\,\theta^jP$, so the array $\mathbf b$, its row sums
$(0,\sigma,\dots,\sigma)$, and $\tau(\mathbf b)$ are unchanged; and
$\mathrm{ord}_{x=0}h_j\ge j+1$, so the tuple has pairwise distinct
vanishing orders at the origin.
\end{lemma}

\begin{proof}
The $n$-th coefficient of $\prod_{i=1}^{j}(\theta-i)P$ is
$(n-1)(n-2)\cdots(n-j)\,c_n$; at $n=0$ this is $(-1)^jj!\,c_0$, and it
vanishes for $1\le n\le j$, whence $\mathrm{ord}\,h_j\ge j+1$ after the
constant subtraction. The operator $\prod_{i=1}^{j}(\theta-i)$ expands
monically in $\Z[\theta]$ (integer Stirling coefficients), and
$cq\,P(0)=cq\,c_0\in\Z$ by \eqref{eq:type}; hence the change of basis
is integral, triangular, with unit diagonal, and $\Q(x)$-linear
independence persists. For the type: the coefficients of $cq\,h_j$ are
the \emph{integer} multiples $(n-1)\cdots(n-j)\cdot cq\,c_n$ of those
of $cq\,P$ (for $n\ge1$; the constant term is $0$), so integrality
against the same denominators $\lcm{n}^{\sigma}$ persists verbatim.
Likewise, multiplication of any admissible function by a fixed positive
integer preserves the type, and no quantity in Theorem~\ref{thm:cdt}
depends on such constants.
\end{proof}

\begin{remark}
We stress that Theorem~\ref{thm:cdt} imposes no hypothesis on the
local behaviour of the $f_i$ at the origin: the vanishing-filtration
data entering its proof depends only on the $\Q$-span of
$\{x^sf_i\}$, which Lemma~\ref{lem:val}'s unimodular change of basis
leaves untouched. The lemma is included so that the application can be
read with either basis in mind.
\end{remark}

\begin{proof}[Proof of Theorem~\ref{thm:main}]
Apply Theorem~\ref{thm:cdt} to the $m:=k+1$ functions
\[
f_1=1,\quad f_{j+2}=cq\,\theta^jP\;(0\le j\le k-1),
\]
on $\Omega=\C\setminus[x_2,\infty)$, with the array $\mathbf b$ having
first row $0$ and the remaining $k$ rows $(1,\dots,1)$ of length
$\sigma$. The hypotheses hold: linear independence by
Lemma~\ref{lem:indep}; the arithmetic type by Lemma~\ref{lem:types};
holonomicity since each $f_i$ lies in the $d/dx$-stable finite
$\Q(x)$-module of Lemma~\ref{lem:types}; holomorphy on $\Omega$ since
$P$ is holomorphic there and $\theta$ preserves
$\mathcal O(\Omega)$; and $\rho(\Omega,0)=4x_2$ by the Koebe map. The
row sums are $(0,\sigma,\dots,\sigma)$, whence
\[
\tau(\mathbf b)=\frac{\sigma}{m^2}\sum_{i=2}^{m}(2i-1)
=\sigma\,\frac{m^2-1}{m^2},
\]
and \eqref{eq:criterion} gives
$\log\rho=\log(4x_2)>\sigma\frac{m+1}{m}
>\sigma\frac{m^2-1}{m^2}=\tau(\mathbf b)$, so the bound applies and
yields
\[
k+1=m\;\le\;\frac{\log(4x_2)}{\log(4x_2)-\tau(\mathbf b)}.
\]
Rearranged, $\log(4x_2)\le\sigma\frac{m+1}{m}=\sigma\frac{k+2}{k+1}$,
contradicting \eqref{eq:criterion}. Hence no such $p/q$ exists.
\end{proof}

\begin{remark}
For $k=1$ the tuple $\{1,P\}$ requires $P\notin\Q(x)$, which an
irreducible order-$1$ operator does not guarantee; restoring
Zudilin's threshold $\tfrac32\sigma$ in that case requires his
positivity hypotheses. For $k\ge2$, irreducibility replaces positivity
altogether.
\end{remark}

\begin{remark}\label{rem:multivalent}
Using the bivalent map $8x_2(z+z^3)/(1+z)^4$ of
\cite[\S2.11]{CDT} together with the evaluation
\cite[Lemma~2.11.7]{CDT} of its Bost--Charles integral
($\log 8+4G/\pi$, $G$ the Catalan constant) in the multivalent bound
\cite[(2.5.4)]{CDT} improves the thresholds further for large $k$; we
have preferred the clean univalent statement.
\end{remark}

\section{The window, and a specification curve for $\zeta(5)$}
\label{sec:window}

Theorem~\ref{thm:main} concludes in the region
\[
\frac{e^{\sigma(k+2)/(k+1)}}{4}\;<\;x_2\;\le\;e^{\sigma},
\]
where the classical criterion ($x_2>e^{\sigma}$) is silent, and which
for every $k\ge2$ strictly contains the determinantal window. (Above
$e^{\sigma}$ the classical argument applies; in fact no transcendental
function of exact type \eqref{eq:type} can have radius of convergence
exceeding $e^{\sigma}$, so the window is also the natural habitat of
such functions.)

For linear forms in $1$ and $\zeta(5)$ one expects $\sigma=5$. A
hypothetical pair $(A,V)$ annihilated by an irreducible operator of
order $k$, with $\zeta(5)A-V$ holomorphic on
$\C\setminus[x_2,\infty)$, proves the irrationality of $\zeta(5)$
provided $4x_2>e^{5(k+2)/(k+1)}$:
\begin{center}
\begin{tabular}{lcccccc}
\toprule
$k$ & $2$ & $3$ & $4$ & $6$ & $10$ & $\infty$\\
\midrule
$x_2>$ & $197$ & $130$ & $101$ & $75$ & $58$ & $37.1$\\
\bottomrule
\end{tabular}
\end{center}
No construction approaching these specifications is currently known;
the table is offered as a target. We note that mixed forms (involving
$\zeta(3)$ as well) can be handled by the same method in the framework
of linear independence statements, at the cost of enlarging the tuple.

\section{Experimental observations on denominators}\label{sec:exp}

The denominator term $\tau(\mathbf b)$ of Theorem~\ref{thm:cdt} is a
worst-case bound. For the Ap\'ery $\zeta(3)$ system
($u_n,v_n$ the Ap\'ery sequences, $\sigma=3$, $k=3$) we computed, in
exact arithmetic, the kernel of the vanishing-to-order-$n$ constraints
for the tuples $\{1,V,\theta V\}$ and $\{1,V,\theta V,\theta^2V\}$
(rational surrogates with the same types and the same
derivative-tower correlation structure as the conditional $P$), for
polynomial degrees up to $22$ and vanishing orders $n^*\le 67$, and
measured the minimal denominator of the leading coefficient over the
kernel. The observed rates are
\[
\frac{\log\operatorname{den}(\beta)}{n^*}\;\approx\;1.76\ (m=3),
\qquad 1.89\ (m=4),
\]
against the theoretical caps
$\tau(\mathbf b)\cdot\frac{3\psi(n^*)}{\sigma n^*}\approx 2.5$--$2.6$
at these ranges; moreover the savings are concentrated at small primes
($p\le n^*/4$), the primes in $(n^*/2,n^*]$ appearing with their full
exponent~$3$. This is consistent with the supercongruence phenomena
known for Ap\'ery numbers and suggests that, for specific systems,
provable ``fine'' denominator types substantially below
\eqref{eq:type} may be available, in the spirit anticipated in
\cite[Remarks~6.6.15 and~10.2.3]{CDT}; any such improvement lowers the
thresholds of Theorem~\ref{thm:main} accordingly.

\section*{Provenance and disclosure}
This note records the outcome of an extended exploration conducted as
a human--AI collaboration. The direction of inquiry, the decisions at
every fork, and the responsibility for circulating this note rest with
the human correspondent (who, out of caution, leaves the byline empty
for now); the mathematical drafting, the literature tracing, and the
numerical experiments of \S\ref{sec:exp} were carried out by the AI
system Claude (Anthropic). Every statement and proof above is
elementary modulo \cite{CDT} and has been written so as to be
checkable by hand against that paper. All of the machinery belongs, of
course, to the authors of \cite{CDT} and \cite{Zudilin}; the present
note merely turns a key that their work cut. Comments, corrections,
and references to prior occurrences of Theorem~\ref{thm:main} are very
welcome and will be gratefully incorporated or, if warranted, the note
withdrawn.

\begin{thebibliography}{9}

\bibitem{Apery} R.~Ap\'ery,
\emph{Irrationalit\'e de $\zeta(2)$ et $\zeta(3)$},
Ast\'erisque \textbf{61} (1979), 11--13.

\bibitem{CDT} F.~Calegari, V.~Dimitrov, Y.~Tang,
\emph{The linear independence of $1$, $\zeta(2)$, and $L(2,\chi_{-3})$},
preprint, \texttt{arXiv:2408.15403}.

\bibitem{vdP} A.~van der Poorten,
\emph{A proof that Euler missed\dots\ Ap\'ery's proof of the
irrationality of $\zeta(3)$},
Math. Intelligencer \textbf{1} (1979), 195--203.

\bibitem{Zudilin} W.~Zudilin,
\emph{A determinantal approach to irrationality},
Constr. Approx. \textbf{45} (2017), 301--310;
\texttt{arXiv:1507.05697}.

\end{thebibliography}

\end{document}
